\section{A uniform quantitative range estimate}
\label{k:section-range}

We now compare the ideal transform with a physical attraction of finite
range.  The estimate is uniform over all preceding attraction ranges and
over arbitrarily small positive added masses.  This uniformity will allow
the error and the next mass to be chosen by a single explicit schedule.
The quadratic density term is the one underlying the Lebowitz--Penrose
limit \citep{LebowitzPenrose1966,GerardiEtAl1973}; our purpose here is a quantitative
estimate with constants uniform across the complete reference family.

\begin{theorem}[Quantitative approximation by a long-range attraction]
\label{k:range}
Fix an admissible $t$ and a compact set $\mathcal D\subset\Omega$.
There are constants $C_{\mathcal D,t}<\infty$ and
$K_{\mathcal D,t}<\infty$ with the following property.  Let $R_0\ge2$,
and let
\[
 \varphi_0=\psi_t-t\sum_j\alpha_jw_{R_j},
 \qquad 1\le R_j\le R_0,
 \qquad \alpha_j\ge0,
\]
where the sum may be empty, finite, or countably infinite.  If
$a>0$, $\sum_j\alpha_j+a\le A+1$, and $R/R_0\ge K_{\mathcal D,t}$,
then
\begin{equation}
 \sup_{\lambda\in\mathcal D}
 \abs{Q_{\varphi_0-ta w_R}(\lambda)-T_aQ_{\varphi_0}(\lambda)}
 \le C_{\mathcal D,t}\sqrt{\frac{R_0}{R}}.
 \label{k:range-estimate}
\end{equation}
Both constants are independent of $R_0$, the number and values of the
reference ranges, the reference masses, and the positive mass $a$.
\end{theorem}

Use boxes of an intermediate side $s$, with $R_0\le s\le R$.
Deleting the reference interaction between different boxes costs
$O(R_0/s)$ in pressure.  The suprema and infima of the new kernel on
pairs of boxes have row sums which, after multiplication by $s^3$,
approximate its unit integral within $O(s/R)$.  Density bounds will give
pressure errors of the same order.  The suprema will give a product upper
bound; the infima will give a lower bound from independent reference
ensembles in the boxes.  Balancing these two errors explains
the choice $s\asymp\sqrt{R_0R}$.  We begin with the deletion estimate,
which must include both the positive and negative parts of the reference
interaction.

\begin{lemma}[Grand pressure in a finite box]
\label{k:box-pressure}
Let $\varphi_0$ and $R_0$ be as in \Cref{k:range}, allowing
$\sum_j\alpha_j\le A+1$.  Write
\[
 Q^s_{\varphi_0}(\lambda)
 =\frac1{ts^3}\log\Xi_s^{\varphi_0}(\lambda).
\]
For every compact $\mathcal K\subset\Omega$ and every $s\ge R_0$ with
$s/b\in\N$,
\begin{equation}
 \sup_{\lambda\in\mathcal K}
 \abs{Q^s_{\varphi_0}(\lambda)-Q_{\varphi_0}(\lambda)}
 \le C_{\mathcal K,t}\frac{R_0}{s}.
 \label{k:box-pressure-bound}
\end{equation}
More generally, for the same $s$, partition $\Lambda_{ms}$ into $m^3$
boxes of side $s$, aligned with the small-cell grid.  In the potential
$\varphi_0-ta w_R$, where $a\ge0$, $R\ge1$, and
$\sum_j\alpha_j+a\le A+1$, delete only reference interactions between
distinct boxes.  If $\Xi^{\mathrm{del}}_{ms}$ is the resulting grand
partition function, then
\begin{equation}
 \sup_{\lambda\in\mathcal K}
 \frac{\abs{\log\Xi_{ms}^{\varphi_0-ta w_R}(\lambda)
                 -\log\Xi^{\mathrm{del}}_{ms}(\lambda)}}{t(ms)^3}
 \le C_{\mathcal K,t}\frac{R_0}{s}.
 \label{k:deletion}
\end{equation}
The constants have the uniformity asserted in \Cref{k:range} and are
also independent of $m,s,R$.
\end{lemma}

\begin{proof}
Put $S_0=t\sum_j\alpha_jw_{R_j}$ and $S_{\mathrm{new}}=ta w_R$.
The series for $S_0$ converges uniformly, is supported in the ball of
radius $R_0$, and has the row bound \eqref{f:row-bound}, including when
there are countably many reference terms.
Multiply the reference interaction on pairs belonging to distinct large
boxes by a parameter $\theta\in[0,1]$.  Every intermediate interaction
retains all intra-small-cell core terms, and its attractive part is
bounded by $S_0+S_{\mathrm{new}}$.  This has small-cell row sum at most
$M$.  Therefore \Cref{f:moments,f:perturbation} apply throughout this
interpolation, whether or not the intermediate interaction is translation
invariant.

The absolute difference is bounded on every pair by
\[
 (\psi_t+S_0)(x-y)
 \boldsymbol 1\{x,y\text{ belong to distinct large boxes}\}.
\]
Its row sum over small cells is at most
$C_b t^2+c_b t(A+1)$: the core is bounded by $t^2$ and supported within
distance two, and the attraction has the common row bound already proved.
This estimate uses an absolute majorant and thus remains valid when
$\varphi_0$ changes sign on a pair of cells.

A nonzero difference row must belong to a small cell within distance
$R_0+\sqrt3 b$ of the boundary of its large box.  Since $R_0\ge2$ and
$s\ge R_0$, there are at most $C_b s^2R_0$ such small cells per large
box.  The bound follows either by subtracting the volume of the inner
cube when the strip is thin, or from $s^3\le C s^2R_0$ when the strip
fills a fixed fraction of the box.  Thus the sum of all difference
majorants is at most
\[
 C_b m^3s^2R_0\bigl(C_b t^2+c_b t(A+1)\bigr).
\]
The finite-volume interpolation bound of \Cref{f:perturbation}, divided
by $t(ms)^3$, proves \eqref{k:deletion}.

For $a=0$ the deleted system consists of independent reference boxes,
and its grand partition function is exactly
\begin{equation}
 \Xi^{\mathrm{del}}_{ms}
       =\bigl(\Xi_s^{\varphi_0}\bigr)^{m^3}.
 \label{k:factorization}
\end{equation}
For completeness, with prescribed box counts $(n_z)$ of total $N$, the
$N!/\prod_z n_z!$ assignments of labels to boxes cancel the global
$1/N!$ and leave $\prod_z1/n_z!$.  Summing the nonnegative product
integrals over all $(n_z)$ gives \eqref{k:factorization}, including
the zero-count terms.  Its normalized logarithm equals
$Q^s_{\varphi_0}$ for every $m$.  Letting $m\to\infty$ in
\eqref{k:deletion} and using only the pressure limit of the full
reference proves \eqref{k:box-pressure-bound}.
\end{proof}

The next estimate uses the global $C^1$ bound of $w$.  Radial monotonicity
and strict positivity on its support are not required.

\begin{lemma}[Upper and lower kernel row sums]
\label{k:kernel-rows}
Let $B_z$ be the box of side $s$ centered at $z$, where $z$ ranges
over a translated lattice $s\mathbb Z^3$, and suppose $0<s\le R$.
Define
\[
 K^+_{zz'}=\sup_{x\in B_z,\,y\in B_{z'}}w_R(x-y),
 \qquad
 K^-_{zz'}=\inf_{x\in B_z,\,y\in B_{z'}}w_R(x-y).
\]
There is a constant $C_w$ depending only on $w$ and dimension such that,
with $\eta=C_ws/R$, every finite cubical collection of these boxes
satisfies
\begin{equation}
 s^3\sum_{z'}K^+_{zz'}\le1+\eta.
 \label{k:upper-row}
\end{equation}
If $z$ has distance greater than $R+4s$ from the boundary of the finite
collection, then
\begin{equation}
 s^3\sum_{z'\ne z}K^-_{zz'}\ge1-\eta.
 \label{k:lower-row}
\end{equation}
Both matrices are symmetric and nonnegative.
\end{lemma}

\begin{proof}
Symmetry follows from $w_R(x-y)=w_R(y-x)$.  Let
$G=\norm{\nabla w}_\infty$, so $\Lip(w_R)\le GR^{-4}$.
For a fixed $y\in B_{z'}$, each displacement $x-y'$ with
$x\in B_z$, $y'\in B_{z'}$ differs from $z-y$ by at most
$3s$.  Consequently
\[
 \abs{K^\pm_{zz'}-w_R(z-y)}\le3GsR^{-4}.
\]
On the infinite grid, at most $C(R/s)^3$ second boxes can have either
term nonzero: all others lie outside the support enlarged by a distance
$3s$.  Integrating over $y\in B_{z'}$ and summing thus gives
\[
 \abs{s^3\sum_{z'}K^\pm_{zz'}-
                  \int_{\R^3}w_R(z-y)\dd y}
 \le Cs^3(R/s)^3sR^{-4}=C\frac sR.
\]
The integral is one.  This comparison includes boxes meeting internal
zero regions of $w$ or crossing its support boundary, because the same
global Lipschitz estimate applies there.

Restricting the upper sum to a finite collection only decreases it.
For a center with the stated boundary distance, all potentially
contributing second boxes are still present in the lower sum.  Omitting
$z'=z$ loses at most
$s^3K^-_{zz}\le\norm w_\infty(s/R)^3$, which is bounded by
$\norm w_\infty s/R$.  Increasing the one constant $C_w$ proves both
claims.
\end{proof}

\begin{proof}[Proof of \Cref{k:range}]
We first bound the particle counts needed for the upper estimate.
For $(\beta,h)\in\mathcal D$, $s\ge1$, $V=s^3$, and $n\ge0$, put
\[
 B_n=\ee^{thn}Z^{\varphi_0}_{s,n}(\beta)
          \exp\!\left(\frac{\beta ta(1+\eta)}{2V}n^2\right),
 \qquad 0\le\eta\le1.
\]
The reference $\varphi_0$ belongs to \eqref{f:family}, and the
auxiliary quadratic coefficient satisfies $a(1+\eta)\le2(A+1)$.
Apply \Cref{f:grand-tail} on $\mathcal D$ with
$\zeta=a(1+\eta)$.  It supplies a common cutoff
$D_{\mathrm{cut}}\ge1$ and, after increasing a volume threshold
$V_0\ge1$, the bound
\begin{equation}
 \sum_{n>D_{\mathrm{cut}}V}B_n\le1
 \qquad(V\ge V_0).
 \label{k:grand-tail}
\end{equation}
Both $D_{\mathrm{cut}}$ and $V_0$ depend only on $\mathcal D,t$.
They are fixed before any local occupation constants: the tail lemma
uses only the canonical quadratic bound, even with this auxiliary
weight.

Let $U$ be the common active-shift bound on $\mathcal D$ from
\Cref{k:transform-properties}.  Choose one compact rectangle
$\mathcal D'\subset\Omega$ containing $\mathcal D$ and every point
\begin{equation}
 (\beta,h+\beta av),\qquad
 (\beta,h)\in\mathcal D,\quad
 0\le a\le A+1,\quad 0\le v\le\max\{U,D_{\mathrm{cut}}\}.
 \label{k:shifted-compact}
\end{equation}
Only now fix the moment and deletion constants of
\Cref{f:moments,k:box-pressure} on $\mathcal D'$.  All these constants
are independent of the eventual box side and of the reference.

Put $q=R/R_0$ and choose
\begin{equation}
 s=b\left\lfloor\frac{\sqrt{R_0R}}b\right\rfloor.
 \label{k:box-side}
\end{equation}
A single lower bound on $q$, uniform for $R_0\ge2$, ensures
\begin{equation}
 R_0\le s\le R,\qquad
 \frac12R_0\sqrt q\le s\le R_0\sqrt q,\qquad
 s^3\ge V_0,\qquad
 \eta=C_w\frac sR\le1.
 \label{k:ratio-conditions}
\end{equation}
Partition $\Lambda_L$, $L=ms$, into $m^3$ boxes $B_z$, and delete
reference interactions between different boxes as in
\Cref{k:box-pressure}.  The new attraction remains intact.  All limits
in $m$ below hold with $s,R,R_0$, the reference, and the parameters fixed.

\medskip\noindent\emph{Upper bound.}
Let $N_z$ be the particle count in $B_z$, and write
$W_R=\sum_{i<j}w_R(x_i-x_j)$.  Adding nonnegative diagonal self terms,
then using symmetry and $2N_zN_{z'}\le N_z^2+N_{z'}^2$, gives
\begin{align*}
 W_R
 &\le\frac12\sum_{z,z'}K^+_{zz'}N_zN_{z'}
 \le\frac12\sum_zN_z^2\sum_{z'}K^+_{zz'}
 \le\frac{1+\eta}{2s^3}\sum_zN_z^2.
\end{align*}
Splitting labels among boxes as in \eqref{k:factorization}, we obtain
\begin{equation}
 \Xi_L^{\mathrm{del}}(\beta,h)
 \le\left[\sum_{n\ge0}
 \ee^{thn}Z^{\varphi_0}_{s,n}(\beta)
 \exp\!\left(\frac{\beta ta(1+\eta)}{2s^3}n^2\right)
 \right]^{m^3}.
 \label{k:upper-product}
\end{equation}
For $0\le n\le D_{\mathrm{cut}}s^3$, let $d=n/s^3$.  The canonical term at count
$n$ is one of the nonnegative terms in the reference grand sum at the
shifted field $h+\beta ad$.  Therefore
\begin{align*}
 \frac1{ts^3}\log B_n
 &\le Q^s_{\varphi_0}(\beta,h+\beta ad)
                     -\frac{\beta a}{2}d^2
                     +\frac{\beta a\eta}{2}d^2\\
 &\le Q_{\varphi_0}(\beta,h+\beta ad)
                     -\frac{\beta a}{2}d^2
                     +C_{\mathcal D,t}\left(\frac{R_0}{s}+\eta\right)\\
 &\le T_aQ_{\varphi_0}(\beta,h)
                     +C_{\mathcal D,t}\left(\frac{R_0}{s}+\eta\right).
\end{align*}
All shifted parameters lie in $\mathcal D'$ by
\eqref{k:shifted-compact}, so the finite-box comparison applies; the bound
$d\le D_{\mathrm{cut}}$ controls the extra quadratic.  The argument includes $n=0$,
whose term equals one.  There are $\lfloor D_{\mathrm{cut}}s^3\rfloor+1$ such
counts; the remaining terms have total at most one by
\eqref{k:grand-tail}.  Since $T_aQ_{\varphi_0}\ge0$, taking logarithms
in \eqref{k:upper-product} yields
\begin{align}
 \frac1{tL^3}\log\Xi_L^{\mathrm{del}}
 \le T_aQ_{\varphi_0}
 +C_{\mathcal D,t}\left(\frac{R_0}{s}+\eta\right)
 +\frac{\log(\lfloor D_{\mathrm{cut}}s^3\rfloor+2)}{ts^3}.
 \label{k:upper-final}
\end{align}
Combine this with \eqref{k:deletion} and let $m\to\infty$ in the
full physical pressure.  This proves the desired upper estimate with
the three displayed errors.

\medskip\noindent\emph{Lower bound.}
Fix an active shift $u\in[0,U]$ for $T_aQ_{\varphi_0}(\beta,h)$.
Let $\Pi$ be the product of the $m^3$ reference grand measures in the
boxes, all at $(\beta,h+\beta au)$, and put
\[
 d_s=s^{-3}\mathbb E_\Pi N_z.
\]
The common small-cell moment bound on $\mathcal D'$ gives a uniform
bound $d_s\le C_{\mathcal D,t}$.  Indeed each box consists of exactly
$(s/b)^3$ small cells, and Cauchy--Schwarz bounds each mean count by the
square root of its uniformly bounded second moment.

The change of field and insertion of the new attraction give the exact
identity
\begin{equation}
 \Xi_L^{\mathrm{del}}(\beta,h)
 =\bigl[\Xi_s^{\varphi_0}(\beta,h+\beta au)\bigr]^{m^3}
       \mathbb E_\Pi
       \exp\{\beta ta(W_R-uN)\},
 \label{k:jensen-identity}
\end{equation}
where $N=\sum_zN_z$.  The box factorials in $\Pi$ are precisely those
obtained by splitting the full grand measure among boxes.  All terms
are finite by the common quadratic stability bound.  Jensen's inequality
therefore gives
\begin{equation}
 \frac1{tL^3}\log\Xi_L^{\mathrm{del}}
 \ge Q^s_{\varphi_0}(\beta,h+\beta au)
       -\beta au d_s+\frac{\beta a}{L^3}\mathbb E_\Pi W_R.
 \label{k:jensen-lower}
\end{equation}

Discard the nonnegative contributions to $W_R$ from pairs in one box.
For the remaining pairs the kernel infima give
\[
 W_R\ge\frac12\sum_{z\ne z'}K^-_{zz'}N_zN_{z'}.
\]
Entire distinct boxes are independent under $\Pi$, so for $z\ne z'$,
$\mathbb E_\Pi[N_zN_{z'}]=s^6d_s^2$.  No independence inside a box is
used.  Let $m_{\mathrm{int}}$ be the number of centers with distance
greater than $R+4s$ from the outer boundary.  Nonnegativity of the other
rows and \Cref{k:kernel-rows} imply
\[
 \frac1{L^3}\mathbb E_\Pi W_R
 \ge\frac{m_{\mathrm{int}}}{m^3}\,
                    \frac{1-\eta}{2}d_s^2.
\]
For fixed $s,R$, $m_{\mathrm{int}}/m^3\to1$ as $m\to\infty$.
Consequently \eqref{k:jensen-lower} and \eqref{k:deletion}, using the
known limit of the full physical model, show that
\begin{align*}
 Q_{\varphi_0-ta w_R}(\beta,h)
 &\ge Q^s_{\varphi_0}(\beta,h+\beta au)
       +\beta a\left(-u d_s+\frac{1-\eta}{2}d_s^2\right)
       -C_{\mathcal D,t}\frac{R_0}{s}.
\end{align*}
The exact square identity
\begin{equation}
 -u d_s+\frac{1-\eta}{2}d_s^2
 =-\frac{u^2}{2}+\frac{(d_s-u)^2}{2}-\frac\eta2d_s^2
 \ge-\frac{u^2}{2}-C_{\mathcal D,t}\eta
 \label{k:mean-square}
\end{equation}
and \eqref{k:box-pressure-bound} at the shifted parameter now give
\begin{align}
 Q_{\varphi_0-ta w_R}(\beta,h)
 &\ge Q_{\varphi_0}(\beta,h+\beta au)-\frac{\beta a}{2}u^2
       -C_{\mathcal D,t}\left(\frac{R_0}{s}+\eta\right)\notag\\
 &=T_aQ_{\varphi_0}(\beta,h)
       -C_{\mathcal D,t}\left(\frac{R_0}{s}+\eta\right).
 \label{k:lower-final}
\end{align}
In particular no convergence of $d_s$ to $u$, or convergence of $d_s$
at all, is required.  This proves the same bound at reference coexistence
parameters.

Finally, \eqref{k:ratio-conditions} gives
$R_0/s\le2q^{-1/2}$ and $s/R\le q^{-1/2}$.  For $s\ge1$,
\[
 \frac{\log(\lfloor D_{\mathrm{cut}}s^3\rfloor+2)}{ts^3}
 \le\frac{C_{\mathcal D,t}}s
 \le C_{\mathcal D,t}q^{-1/2},
\]
where the last step uses $R_0\ge2$ and
$s\ge\tfrac12R_0\sqrt q$.  Equations
\eqref{k:upper-final} and \eqref{k:lower-final}, with the deletion cost,
prove \eqref{k:range-estimate}.

No step imposes a positive lower bound on $a$, and no error has been
divided by $a$.  Only the reference and full physical models require
thermodynamic limits; the deleted systems were treated by exact
factorization or finite-volume inequalities.
\end{proof}
